\section{Introduction}

\acp{bss} are increasingly relevant in urban mobility as a sustainable and active transport mode~\cite{Fishman2016, MeddinWorldMap2025}. As systems are deployed or expanded, a central planning question is where to place a limited number of stations so that network performance aligns with specific policy priorities. This paper focuses on this strategic problem framing it as follows. Given a set of feasible candidate locations, the goal is to select a fixed number of stations that maximize total system utility according to a planning objective. Rather than relying solely on trip or origin-destination estimates, this study focuses on providing a tool to support strategic station placement under multiple planning perspectives.

Existing research includes location-allocation models~\cite{Frade2015, Mete2018}, \ac{mcdm} methods~\cite{Kabak2018, Bahadori2022}, and exact~\cite{Cintrano2018} or metaheuristic~\cite{Qian2022, Caggiani2020} optimization techniques. However, several limitations remain. First, most models are designed for a single objective~\cite{Kabak2018, Bahadori2022, Fazio2021} or a small predefined set of objectives~\cite{Conrow2018, Qian2022, Caggiani2020, Duran-Rodas2021, Fan2024, Nikiforiadis2021}, making it difficult to adapt them to evolving planning priorities without reformulation. Second, while station placement depends on both the local suitability of candidate sites and the structure of the resulting network, these two dimensions are rarely integrated explicitly within the objective function. Third, many studies rely on coarse spatial representations or simplified distance metrics~\cite{Frade2015, Fazio2021, Veillette2018, Nikiforiadis2021} that do not fully capture the structure of urban cycling networks, often neglecting directionality~\cite{Cintrano2018, Xin2013} and topography~\cite{Fazio2021, Kabak2018, Ebrahimi2022, Tera2023}.

To address these limitations, we develop \textit{FlexBSS}, a flexible framework for BSS station placement with four components: (i) a modular \ac{mcdm} scoring scheme that evaluates candidate sites using configurable spatial factors and adjustable weights, enabling different planning objectives to be represented within a unified formulation; (ii) the integration of network-level metrics to explicitly balance station proximity and system-wide accessibility on directed bicycle networks; (iii) an optional topography module that adjusts distances to account for slope-related cycling effort; and (iv) support for both new deployments and the expansion or reconfiguration of existing systems. The resulting binary combinatorial optimization problem is solved using a \ac{ga}, enabling efficient exploration of large candidate sets under diverse planning scenarios.

Conceptually, the framework is intended to support a shift from fixed, single-purpose station-location optimization--where a single objective or a narrow set of objectives is predefined--toward an adaptive, scenario-based decision support tool. Within this same unified formulation, planners explore different planning objectives and compare the resulting station networks and trade-offs. We demonstrate the approach in Barcelona using open data, evaluating scenarios focused on demand, multimodal integration, and social equity, as well as an expansion case in l'Eixample district. Results show how the framework helps planners explore alternative configurations, assess trade-offs, and design context-sensitive station networks.

\section{Literature review} \label{sec:lit_review}

\ac{bss} station placement studies differ in data inputs, modelling approaches, and planning objectives. This review synthesizes the main methodological strands and identifies the gaps addressed in this study.


\subsection{From heuristics to computational planning}

Docked \ac{bss} planning is typically an iterative process that begins with clearly defining the planning goals, such as maximizing ridership, promoting social equity, or improving first-last kilometer access~\cite{Duran-Rodas2021, Ferrando2007, Yanocha2018}. These goals should be used to guide each stage of site selection and refinement~\cite{Ferrando2007, Yanocha2018}, drawing on demand patterns, system performance measures, and stakeholder feedback. In practice, early \ac{bss} deployments relied on expert judgment and heuristics, with planning guides recommending stations in visible, accessible, and well-connected high-demand areas and encouraging stakeholder input~\cite{Ferrando2007, Yanocha2018, CROW2006}. Furthermore, cycling infrastructure manuals emphasize coherent, continuous cycling networks~\cite{AASHTO1999bicycle, CROW2006}. While these guidelines provide valuable practical recommendations, they lack analytical tools or optimization frameworks for systematically evaluating alternatives across diverse urban conditions. Nonetheless, heuristics and optimization can complement each other: practical rules can inform modelling choices, while optimization outputs can help refine or validate heuristic guidelines. As the availability of detailed spatial data and advanced computational tools increased in the early 2010s, researchers moved beyond purely heuristic approaches and began applying formal spatial optimization models.


\subsection{Location-allocation models} \label{sec:location_allocation_models}

Location-allocation models are spatial optimization tools that place a fixed number of facilities and allocate demand to them to optimize one performance metric~\cite{Daskin1995}. Demand is typically defined over zones or grid cells and facilities over candidate sites, usually as discrete-location problems~\cite{Teixeira2008}. Core formulations include the \textit{p}-median problem, which minimizes total demand-weighted distance~\cite{Hakimi1965}; the \ac{mclp}, which maximizes demand served within a service radius (i.e., a threshold within which users are considered covered)~\cite{Church1974}; and the \textit{p}-center problem, which minimizes the maximum distance between users and their nearest facility~\cite{Hakimi1964}.

Several studies have applied these formulations to \ac{bss} placement~\cite{Frade2015, Cintrano2018, Lin2011, Mete2018}. In these applications, models are typically single-objective, most often maximizing coverage or minimizing user distance while incorporating practical constraints. For example, Cintrano et al.~\cite{Cintrano2018} applied a \textit{p}-median formulation; Frade and Ribeiro~\cite{Frade2015} used an \ac{mclp} variant to maximize population coverage under budget constraints while optimizing station capacity, locations, and fleet size; Lin and Yang~\cite{Lin2011} proposed a mixed-integer nonlinear model minimizing system costs (travel, installation, inventory, and unmet demand); and Mete et al.~\cite{Mete2018} evaluated several coverage-based formulations (\textit{p}-center, \textit{p}-median, set covering).

While these models offer reproducible alternatives to heuristic siting, they remain focused on geographical coverage and user-to-station distance, and broader priorities such as social equity or multimodal integration are largely absent.


\subsection{Combining multiple data sources} \label{sec:combining_multiple_data_sources}

Many studies combine multiple data sources to pursue diverse planning objectives and therefore face the challenge of how to aggregate or balance them. The literature generally structures this in two ways. In \textit{multi-criteria} approaches, multiple data sources are combined to support one planning objective; in \textit{multi-objective} approaches, factors associated with different goals are kept separate and trade-offs are made explicit.

In \textit{multi-criteria} studies, the aim is a single planning objective--usually maximizing demand coverage--and the information used to pursue it comes from several sources. Variables such as population density, proximity to \ac{pt}, and cycling infrastructure are typically weighted and aggregated to support that objective. Because these are in different units and play different roles, studies must choose how to weight and combine them. Some studies rely on structured judgment: \ac{ahp}, Fuzzy \ac{ahp}, \ac{moora}, \ac{bwm}, or \ac{topsis} are used to prioritize criteria with stakeholders~\cite{Kabak2018, Guler2021a, Guler2021b, Bahadori2022}. Others derive weights from data using statistical models of observed \ac{bss} usage (e.g.\ regression coefficients or variable importance scores in \ac{ml})~\cite{Wang2021}. Others use either equal~\cite{Fazio2021, Tera2023} or flexible~\cite{Veillette2018} weighting schemes, emphasizing that these should be tailored to local priorities and consulted with stakeholders.

This logic appears in two main forms. In the first, suitability maps are built without formal optimization to rank candidate station locations. For instance, Fazio et al.~\cite{Fazio2021} aggregate thematic indexes into a composite score; Kabak et al.~\cite{Kabak2018} use \ac{ahp} and \ac{moora}; and Tera et al.~\cite{Tera2023} use MULTIMOORA. In the second form, the same aggregation logic is embedded in an optimization model, where multiple spatial factors are combined into a single objective optimized within location-allocation or related frameworks. Some examples include García-Palomares et al.~\cite{GarciaPalomares2012}, which combine building-level and zonal data in a slope-aware network to estimate demand, applying $p$-median and \ac{mclp}; Chou et al.~\cite{Chou2019}, which combine \ac{pt} data in a network flow model with proportional flow constraints to align station locations and dock capacity with latent short-distance demand applying a mixed integer programming model; and Ebrahimi et al.~\cite{Ebrahimi2022}, which use \ac{mclp} and Target Market Share models to evaluate coverage and identify accessibility gaps.

\textit{Multi-objective} approaches, by contrast, treat different planning goals as distinct. No single solution optimizes all objectives at once; instead, the outcome is a set of Pareto-optimal solutions. This separation makes the contribution of each objective explicit and enables systematic exploration of alternatives, rather than concealing them within a single aggregated metric~\cite{Rahimi2023, honegumi}. 

Studies differ in the objectives selected. Conrow et al.~\cite{Conrow2018} maximize coverage of both residential population and bicycle network infrastructure, but a recurring theme is social equity, with several studies balancing demand with equity. Duran-Rodas et al.~\cite{Duran-Rodas2021} generate alternative plans under different demand-equity emphases; Caggiani et al.~\cite{Caggiani2020} minimize accessibility disparities across population groups via a Theil index; Fan and Harper~\cite{Fan2024} jointly maximize demand coverage and \ac{pt} access for underserved communities; and Qian et al.~\cite{Qian2022} optimize station reallocation to improve both revenue and equitable access to jobs and essential services. Multi-objective approaches have also been used for operational goals, such as Nikiforiadis et al.~\cite{Nikiforiadis2021}, who apply a two-phase framework in which objectives are first optimized independently and then balanced through weighted deviations from target values.

Together, these strands show how multiple data sources and goals can be integrated in \ac{bss} planning. In both, however, factor and objective sets are often fixed in advance, so introducing new dimensions typically requires model reformulation. To address this limitation, we introduce a modular \ac{mcdm} framework that allows new criteria to be incorporated without changing the overall structure, and flexible weights so that different planning objectives can be emphasized as needed.


\subsection{Spatial granularity and network representation}

Studies vary in the spatial scope of the planning area and in the granularity of candidate site definition. In terms of scope, analyses range from entire cities~\cite{Conrow2018, Duran-Rodas2021} to districts~\cite{GarciaPalomares2012} or campuses~\cite{Mete2018}. Common granularity choices include zonal formulations~\cite{Frade2015, Mete2018, Conrow2018}, grid-based studies~\cite{Fazio2021, Tera2023, Veillette2018} (e.g.\ 100\,m or 300\,m grids), hybrid combinations of grid cells with zones and transport nodes~\cite{Nikiforiadis2021}, and Voronoi-based catchment zones~\cite{Duran-Rodas2021}. These choices improve tractability but reduce spatial precision. A second group uses network-based approaches, selecting candidates along major corridors, at street intersections, or near \acp{poi} such as \ac{pt} stops. Conrow et al.~\cite{Conrow2018}, Caggiani et al.~\cite{Caggiani2020}, and Fan and Harper~\cite{Fan2024} place candidate sites along the bicycle or \ac{pt} network, offering more spatial detail than coarse zoning but without enumerating all possible locations. Lastly, high-resolution approaches based on street-network-derived candidate sets remain less common: Cintrano et al.~\cite{Cintrano2018} use 33,550 street segments in Málaga, while Xin~\cite{Xin2013} evaluates 1,489 street intersections in downtown Vancouver.


Beyond site definition, some studies use network representations to evaluate travel paths rather than euclidean or zonal approximations~\cite{GarciaPalomares2012, Cintrano2018, Xin2013}, improving on simpler models. Networks are typically represented as undirected, with links treated as bidirectional and of equal length in both directions. This simplification omits one-way streets and asymmetric accessibility but still captures cycling conditions better than non-network models.

Additionally, slope strongly influences \ac{bss} cycling behaviour, with fewer trips observed on uphill trips and more on downhill ones~\cite{Morency2015, MateoBabiano2016}, especially in cities with heterogeneous topography, such as Barcelona~\cite{Grau-Escolano2024} and Brisbane~\cite{MateoBabiano2016}. In terms of modelling, García-Palomares et al.~\cite{GarciaPalomares2012} took an important step by incorporating slope as a uniform penalty on effective travel time, but their approach does not distinguish between uphill and downhill travel, and therefore does not differentiate between the increased effort required for going uphill, the facilitating effect of moderate declines and the safety-related penalties of steep descents.

Building on prior high-resolution~\cite{Cintrano2018, Xin2013} and slope-aware work~\cite{GarciaPalomares2012}, we address two limitations: the use of undirected representations and uniform slope penalties. We employ a directed street-level network to capture circulation constraints (e.g., one-way streets), where shortest paths between two points can differ by direction and may exhibit asymmetries above 10\%~\cite{Melo2022}. In addition, we replace the uniform slope penalty with an effort-based measure that differentiates uphill and downhill gradients, consistent with evidence that uphill speed reductions are near-linear whereas downhill effects are non-linear due to safety-related braking on steeper descents~\cite{Ryeng2016, Flugel2019}. Together, these improvements increase model realism.


\subsection{Spatial arrangement of stations}

An important \ac{bss} design aspect is the spatial arrangement of stations, that is, how stations are distributed throughout the service area. Classical network-based models, such as the $p$-median or $p$-center problems~\cite{Hakimi1964, Hakimi1965}, are typically formulated to enhance accessibility. In these approaches, facility arrangement emerges as a side effect of accessibility optimization rather than as a directly controllable planning goal. This pattern also holds in \ac{bss} applications~\cite{Frade2015, Cintrano2018, Lin2011, Mete2018}, where, despite implementation differences, optimization primarily targets accessibility and treats station spatial distribution as a byproduct rather than an explicit objective that can be controlled.

Beyond accessibility, a complementary dimension of spatial arrangement is proximity, reflecting how close stations are to one another. Accessibility and proximity are distinct dimensions: a tightly clustered set of stations at one edge of the service area may have high inter-station proximity but low accessibility for most users, whereas relocating the same cluster to a central area leaves proximity unchanged while improving accessibility. Proximity can support short trips and rebalancing, yet most models impose only minimum inter-station distances rather than explicitly optimizing proximity~\cite{Nikiforiadis2021, Kabak2018}. Network-science concepts such as farness~\cite{Sabidussi1966}--the sum of shortest-path distances from a node to all others--offer a principled way to quantify spatial separation and complement accessibility metrics.

Building on earlier work emphasizing accessibility or minimum spacing~\cite{Frade2015, Cintrano2018, Lin2011, Mete2018, Nikiforiadis2021, Kabak2018}, we combine accessibility and proximity in a weighted composite score, making spatial arrangement trade-offs explicit within the planning framework.


\subsection{Algorithmic approaches}

To obtain a single optimal station configuration, the station placement problem is usually formulated as a constrained optimization model in which the chosen locations maximize or minimize a defined objective function. Common formulations include the $p$-median, maximal covering, and $p$-center models. These problems are typically set up as integer programmes with binary decision variables indicating whether a station is placed at each candidate location~\cite{Frade2015, Lin2011, Mete2018, Cintrano2018}, and constraints such as budget limits or a fixed number of stations~\cite{Conrow2018, Nikiforiadis2021, Caggiani2020}. They are commonly solved with exact solvers such as CPLEX~\cite{Mete2018, Chou2019}, Gurobi~\cite{Mix2022}, or similar tools~\cite{Lin2011, Conrow2018}, which guarantee optimality but can struggle as the candidate set or problem complexity grows. Some studies therefore use location-allocation tools in \ac{gis} software such as ArcGIS~\cite{GarciaPalomares2012, Ebrahimi2022}, which implement heuristics for $p$-median or covering problems and can quickly generate near-optimal solutions.

To address scalability limits of exact methods, researchers in various fields have increasingly adopted metaheuristic algorithms~\cite{Tomar2023}. Techniques such as \ac{ga}, \ac{simann}, \ac{pso}, and \ac{aco} have been applied to large combinatorial problems in facility location~\cite{Chadry2003, Chalupa2019}, transport planning~\cite{Holliday2023, Sachan2024}, and network design~\cite{Gallo2010, Madadi2024}. Although originally developed for single-objective settings, these methods have been adapted to multi-objective optimization by approaches like weighted aggregation or the $\varepsilon$-constraint method~\cite{Rahimi2023}. However, such adaptations often require multiple algorithm runs to approximate the Pareto front and may fail when trade-offs are non-convex or discontinuous. To overcome these limitations, dedicated multi-objective metaheuristics (e.g.\ \ac{nsgaii}~\cite{Deb2002}, \ac{spea2}~\cite{Zitzler2001} or \ac{mopso}~\cite{Coello2002}) have been developed to approximate the Pareto front in a single run by maintaining a set of nondominated solutions. In \ac{bss}, single-objective metaheuristics have been used to optimize station placement, with \acp{ga} applied in particular to $p$-median and multi-objective problems~\cite{Cintrano2018, Qian2022, Liu2015, Caggiani2020}. Comparative work suggests that \acp{ga} can outperform alternatives such as \ac{simann}, \ac{pso}, or Iterated Local Search on $p$-median formulations~\cite{Cintrano2018}, and hybrid \ac{ga}--\ac{pso} methods have also been explored~\cite{Askari2017, Chen2020}.

Metaheuristic validation typically considers solution quality, computational time, and memory use~\cite{Osaba2021}, and the standard practice is to run multiple independent trials and report summary statistics (e.g., best, mean, standard deviation)~\cite{Osaba2021}. Solution quality is commonly assessed either against known optima or tight bounds~\cite{Roshanaei2010}, or--when exact benchmarks are unavailable--against alternative heuristics/metaheuristics, often with statistical tests to evaluate significance~\cite{Derrac2011, Danka2013}.

In this framework, we adopt a \ac{ga} due to (i) its prior use in \ac{bss} optimization; (ii) its ability to incorporate custom, constraint-aware operators; (iii) its suitability for encoding solutions at the street-network level; and (iv) its stronger performance in preliminary tests. Furthermore, in contrast to most previous \ac{bss} optimization studies, we evaluate the proposed \ac{ga} by benchmarking it against a set of baseline approaches.

